Deployed large language model (LLM) systems face continuously evolving input distributions, yet standard drift detectors rely on low-order statistics that may miss structural reorganization of the embedding manifold. We investigate whether persistent-homology-based drift detectors can detect distribution shift in LLM embedding streams more sensitively than classical methods. In a controlled comparative study using \agnews sentence embeddings across five drift scenarios, we find that classical statistical detectors (centroid shift, covariance shift, MMD) achieve near-perfect AUC\,=\,1.0 on all real-world topic drift scenarios, while TDA methods reach AUC 0.49--0.93. However, a supplementary synthetic experiment reveals where TDA has genuine advantage: when drift changes the \emph{topological structure} (loops, holes) while preserving the centroid, TDA H1 features achieve $12.3\times$ separability versus $0.8\times$ for centroid shift. Our results show that TDA-based monitoring adds value primarily for topology-specific drifts---not for the natural domain shifts commonly tested in benchmarks---and suggest a complementary monitoring strategy that combines fast statistical detectors with TDA features for detecting geometric reorganization invisible to moment-based methods.
